\documentclass[modern]{desc-tex/styles/lsstdescnote}

\usepackage{xcolor}
\newcommand{\eac}[1]{{\em Eric:  }{\color{red}{#1}}}

% LATEX PACKAGES
\usepackage{desc-tex/styles/lsstdesc_macros}

\begin{document}

\title{Photometric Redshift Data Challenge}
\author{E.\ Charles}
\author{Add Yourself} 
\date{\today}

\begin{abstract}
 Addressing DESC science questions will require running analysis
 pipelines on simulated data to ensure these pipelines work as
 intended, characterize their performance, and compare them to inform
 decisions on fiducial choices on data.  We are organizing the
 photometric redshift (PZ) data challenge to asses the current status
 of PZ algoritms and analysis pipelines.   This document describes the
 format of the PZ challenge, including the data that participants will
 use, the types of tasks that they will be asked to do, and the metric
 and criteria by which submission will be assseed.
 \end{abstract}

\maketitle

\section{Introduction}

Redshift inference is a key measurement for multiple DESC science
goals and redshift uncertainty is one of the leading contributors to
overall uncertainty on cosmological models from imaging survey
data. Precursor surveys took a variety of approaches to this problem,
accounting for differences in underlying data as well as modeling
approaches. In all cases, redshift uncertainty was significantly
larger than DESC Science Requirements listed in the LSST DESC Science
Requirements Document.

This state-of-the-art motivates a data challenge to characterize and
improve existing methods, as well as provide infrastructure for the
development of improved methods. Overall, this requires generating
uniform input catalogs to use and infrastructure for comparing output
redshift posteriors to each other and to simulated truth catalogs.


\section{Challenge Format}

The PZ data challenge comprises a set of tasks for participants.
These task will be used to evaluate how ready various algorithms are
to be used for cutting edge analysis.  Readiness will be evaluated on
a few different fronts.  1) Does the algorithm meet performance
requirements?  2) Is it robust, flexible and relatively easy to use on
different datasets?  3) Is it scalable up to the scales we will need
to use it at.

This document and the associated web pages describe the data being
provided to participants, the task they will be asked to perform, and
the metrics by which the algorithms readiness will be evaluated.


\subsection{Scope and Timeline}

The data challenge will include two major parts, with a set of
task emulating increasingly realistic scenarios in each part.   The
first part, $p(z)$ estimation, will focus on estimating the redshift of
individual objets.   The second part, tomogrpahy and $n(z)$ estimation
will focus on assigning object to tomographic bins and estimating the
distribution of redshifts in each bin.

The data challenge will run from April 13, 2026 to August 7, 2026.
The first set of data and tasks related to $p(z)$ estimation will be
released on April 13.  A second set of data and tasks related to
$n(z)$ estimation will be released on June 1.

Preliminary results will be released on August 14, 2026, with a
techincal note summarize those result to follow shortly thereafter and
a comprehensive journal publication to follow later.


\subsection{Challgenge Data}

The preparation of the challenge data is described in the appendices.
The data is available as \texttt{tar} archives on the data challenge site.

Each task set in the data challenge has an associated set of files.
Typically these will be a collection of training files that contain
photometric data and reference redsfhits and a second set of files
that contain photometric data but do not include redshifts.   Each
task set will invovle estimating something about the redshifts or
redshift distributions in the test files.

Typically there will be several training and test files for a
particular test set, covering different scenarios and using different
input simulations.


\subsection{Challenge Tasks and Submission Types}

The challenge is organized as a series of sets tasks using
increasingly realistic representations of the data.   In general, each
set of task includes 3 tasks.

\begin{enumerate}
  \item{Estimate either per-object $p(z)$ or ensemble $n(z)$
      distributions for a set of different scenarios and provide the
      estimates in a specfied format.}
  \item{Provide trained models for the different scenarios, and a
      generic script that can be used generate the estimates from task
      1 on an arbitrary dataset.}    
  \item{Provide a generic script that can be used generate the
      generate the models and estimates from tasks 1 and 2 on
      arbitrary datasets.}
\end{enumerate}


\subsection{Input data format}

The input data for the challenge are presented in hdf5 files.
The naming convention for the files is
\texttt{\{challenge\}\_\{taskset\}\_\{simulation\}\_\{label\}\_\{scenario\}.hdf5}.
The meaning for the various fields are descrbied in
Tab.~\ref{tab:file_fields}.   The columns in the files are described
in Tab.~\ref{tab:columns}.

\eac{Do we mention tables\_io for conversion}


\begin{table}[ht]
\centering
\caption{\label{tab:file_fields}\small Fields in .}
\begin{tabular}{ll}
  \hline\hline  
  field & Description \\ \hline
  challenge & Challenge associated to file \\
  taskset & Task set associated to file \\
  simulation & Simulation used to produce file \\
  label & File label (e.g., ``test'', ``training'' \\
  scenario & Data Scenario \\
\hline\hline
\end{tabular}
\end{table}


\begin{table}[ht]
\centering
\caption{\label{tab:columns}\small Contents of input files.}
\begin{tabular}{ll}
  \hline\hline
  Column & Description \\ \hline  
  redshift & True redshift (training files only) \\
  ra & Right Accension (training files only) \\
  dec & Declination (training files only) \\
  object\_id & Unique Object ID\\
  mag\_\{band\}\_lsst & Magnitude in LSST \{band\} \\
  mag\_\{band\}\_lsst\_err & Magnitude uncertainty in LSST \{band\} \\
  mag\_\{band\}\_roman & Magnitude in Roman \{band\} \\
  mag\_\{band\}\_roman\_err & Magnitude uncertainty in Roman \{band\} \\
\hline\hline
\end{tabular}
\end{table}



\subsection{Data format for per-object $p(z)$ estimates}

The $p(z)$ estimates should be submitted as in \texttt{qp} format,
which allows users to specific a complete $p(z)$ distribution for
each object, as well as summary statistics for each object.   

The \texttt{qp} packages supports several different representation of
$p(z)$, such as different functional forms as well as interpolated
grids, histograms, and others.  

For users unfamiliar with \texttt{qp} we highly recommend representing
the $p(z)$ either an interpolated grid, or a Gaussian mixture model.

\begin{verbatim}
# Interpolated grid
import qp
import numpy as np
# Define the x-grid.  Note that we put all the
# p(z) on the same x-grid
xvals = np.array([0,0.5,1,1.5,2])
# Define the y-values.  Note we provide n_grid_points x n_objects 
# values, as we need to provide a y-value at each grid point 
# for each object.
yvals = np.array(
 [
   [0.01,0.2,0.3,0.2,0.01],
   [0.1,0.3,0.5,0.2,0.05]
 ]
)
ens = qp.interp.create_ensemble(xvals,yvals)
\end{verbatim}

\begin{verbatim}
# Mixture model
import qp
import numpy as np
# Define the means, standard deviations and weights.
# These should each have shape n_objectws, n_components.
# In this case we are defining 3 objects with 2-Gaussian 
# representations.
# For each object the weights should sum to 1, or they
# will be normalized.
means = [[0.3, 0.4], [0.5, 0.5], [0.6, 0.8]]
stds = [[0.2, 0.4], [0.1, 0.3], [0.05, 0.3]]
weights = [[0.8, 0.2], [0.7, 0.3], [0.8, 0.2]]
ens = qp.mixmod.create_ensemble(means=means,stds=stds,weights=weights)
\end{verbatim}

The submission files should use the same file name conventions defined
in Tab.~\ref{tab:file_fields}.   The labels will typically be
\texttt{pz\_estimate} or \texttt{pz\_model}, and will be specified in
the descriptions of the various tasks, e.g., 
\texttt{pz\_challenge\_taskset\_1\_cardinal\_pz\_estimate\_yr1.hdf5} or
\texttt{pz\_challenge\_taskset\_1\_cardinal\_pz\_model\_yr1.pkl}.


%\subsection{Data format for tomographic bin assignments}

%\eac{Details about data packaging}


%\subsection{Data format for ensemble $n(z)$ estimates}

%\eac{Details about data packaging}


\subsection{Format for estimation only scripts and trained models}

For the second sub-task, providing trained models and a script to run
estimation using those trained models, submitters should start with
the \texttt{run\_estimate.sh} example script, and modify the
portions of the script to install their packages, download their
models, and run the estimation on the appropriate test files.  

Once they have tested the script, submitters should open a pull
request to add \texttt{run\_estimate.sh}



\subsection{Format for training and estimation scripts}

For the second sub-task, providin a script to train models and run
estimation using those trained models, submitters should start with
the \texttt{run\_train\_and\_estimate.sh}  example script and modify the portions of the script to install
their packages, download their models, and run the estimation on the
appropriate test files.  


\section{Metrics and Assesment Criteria}


\subsection{Metrics for per-object point-estimates}

\eac{Bias, scatter and MAD}


\subsection{Metrics for per-object $p(z)$ distributions}

\eac{PIT-QQ related metrics}


\subsection{Metrics for per-object estimation computational performance}

\eac{Training time, per-object estimation time, model size, estimate data volume per-object}


%\subsection{Metrics for bin assignement}

%\eac{Efficiency and purity}


%\subsection{Metrics for ensemble $n(z)$ estimation}

%\eac{Bias, width, simplified fisher forecast}


%\subsection{Metrics for ensemble estimation performance}

%\eac{Training time, per-object estimation time, model size}



\section{Challenge Tasks related to $p(z)$ estimation}\label{sec:tasks}

\subsection{Task set 1:  estimate redshifts using representative
  training samples}

The first, simplest, task is to estaimate redshifts using
representative training samples.  I.e., the training samples are drawn
for the same distibutions as the test samples.  For this task set we
did not use an of the spectroscopic selection emulation, but simply
applied a uniform magnitude cut $i < 23$ in selecting objects for both
the training and test samples.

The four 
\texttt{pz\_challenge\_taskset\_1\_\{simulation\}\_training\_\{scenario\}.hdf5}
files are the training sets for the ``Flagship'' and ``Cardinal''
simulations, and emulating 1 year and 10 years of LSST data under
expected observing strategy and conditions.  These files have true
redshifts to serve as labels.

The corresponding 
\texttt{pz\_challenge\_taskset\_1\_\{simulation\}\_test\_\{scenario\}.hdf5}
files were drawn from the same distributions.  The true redshifts have
been removed from these files.   The task is to assign $p(z)$
estimates for all the objects in thse 4 test files.



\subsection{Task set 2:  estimate redshifts on non-representative
  samples}

The second, slightly more challenging, task is to estimate redshifts
using non-representative training samples.  I.e., the training samples
are not drawn for the same distibutions as the test samples.  For this task set we
applied the spectroscopic selection emulation for the train set, but
retained all the objects down to $i < 25.4$ in the test set.
Accordingly, the training set will not be representative of the
fainter objects in the test set.   This reflects that spectroscopic redshifts
are typically significantly more diffiicult to obtain than photometry.

The four 
\texttt{pz\_challenge\_taskset\_1\_\{simulation\}\_training\_\{scenario\}.hdf5}
files are the training sets for the ``Flagship'' and ``Cardinal''
simulations, and emulating 1 year and 10 years of LSST data under
expected observing strategy and condition and with spectroscopic
selections emulated.

The corresponding 
\texttt{pz\_challenge\_taskset\_1\_\{simulation\}\_test\_\{scenario\}.hdf5}
files were drawn from the distributions of all objects down to $i <
25.4$, and the true redshifts have been removed from these files. 
The task is to assign $p(z)$ estimates for all the objects in thse 4 test files.




%\subsection{Task set 3:  estimate redshifts on non-representative
 % samples with some mis-labeled redshifts}


%\subsection{Task set 4:  estimate redshifts on non-representative
 % samples with some mis-labeled redshifts and blended objects}



%\section{Challenge Tasks related to $n(z)$ estimation}


%\subsection{Task set 1:  assign objects to fixed tomographic bins and
%  estimate the n(z) distribution in each bin}


%\subsection{Task set 2:  assign objects to arbitrary n(z) bins and
%  estimate the n(z) distribution for each bin}



\appendix

\section{Input simulations}

The challenge employs simulated galaxy catalogs derived from two
complementary N-body cosmological simulations: the Cardinal
simulations and the Flagship simulation. These synthetic datasets
provide a controlled environment where the true redshifts are known by
construction, enabling rigorous validation of photometric redshift
algorithms and systematic assessment of their performance
characteristics.

The Cardinal simulations comprise a suite of high-resolution N-body
simulations specifically designed to explore the sensitivity of
cosmological observables to variations in fundamental cosmological
parameters.   The simulations employ state-of-the-art semi-analytic
models to populate dark matter halos with galaxies, incorporating
realistic prescriptions for star formation, dust attenuation, and
spectral energy distribution modeling.

The Flagship simulation represents a single, ultra-large cosmological
simulation run with fiducial cosmological parameters consistent with
current observational constraints. With a volume exceeding several
cubic gigaparsecs, the Flagship provides statistical power to probe
rare objects and the high-mass end of the galaxy population. Its
primary purpose in the photometric redshift challenge is to provide a
realistic mock catalog that captures the full complexity of galaxy
populations across cosmic time, including correlations between galaxy
properties, environmental dependencies, and the intricate
relationships between spectral features and redshift.

Together, these complementary simulation suites enable challenge
participants to test both the accuracy and the robustness of their
photometric redshift estimation methods under realistic observational
conditions.

\section{Emulating observational effects}

To bridge the gap between the idealized simulation outputs and
realistic survey observations, we employ the RAIL (Redshift Assessment
Infrastructure Layers) software package to emulate observational
effects. RAIL provides a modular framework for injecting realistic
photometric uncertainties, applying survey-specific selection
functions, and simulating the measurement errors characteristic of
modern large-scale imaging surveys. This processing ensures that the
simulated galaxy catalogs reflect the complexities of actual
observations, including magnitude-dependent photometric scatter,
incomplete sky coverage, and the effects of source blending in crowded
fields, thereby providing a more stringent and realistic test bed for
photometric redshift estimation algorithms.

\subsection{Photometric Smearing}

Central to our observational emulation is RAIL's wraping of the
photometric error module, photErr, which we have extended and wrapped
to account for realistic observing strategies and time-dependent
survey conditions. The standard photErr module provides basic
photometric error modeling based on magnitude-dependent noise
characteristics, but our enhanced version incorporates additional
complexity including spatially-varying depth maps. This
wrapper accesses detailed operational simulation outputs that emulate
the expected LSST survey strategy.

Our photErr implementation computes photometric uncertainties by
combining the intrinsic Poisson noise from source photons with
realistic models of sky background, readout noise, and other
systematic contributions.  For each simulated galaxy, use the expected
co-addition depth to derive final photometric error estimates. This
approach captures the heterogeneous nature of survey depth across the
footprint, where some regions benefit from numerous high-quality
exposures while others may be observed only during poor
conditions. The resulting photometric uncertainties vary realistically
with position on the sky, band-dependent limiting magnitudes, and
local observing history, providing challenge participants with mock
catalogs whose noise properties more closely match those expected from
the actual survey.

\subsection{Spectroscopic and narrow-band photo-metric redshift
  selection}

RAIL can emulate the selection functions of several different
spectroscopic redshift surverys, including VVDSf02, zCOSMOS,
DEEP2\_LSST, and the DESI BGS, ELG and LRG samples.

We can also use RAIL to emulate narrow-band photometric surveys,
include small amounts of mis-labeled reference redshifts.


\section{Preparing Training and Test and Reserved datasets}

All of the data prepration was performed using the
\texttt{rail\_projects} and \texttt{rail\_package\_config} packages
for bookkeeping and reproducibility.

\begin{table}[ht]
\centering
\caption{\label{tab:prep_scripts}\small Scripts used
  in data preparation.}
\begin{tabular}{lll}
\hline \hline
Script & Command Run & Purpose \\ \hline
  do\_00\_reduce & rail-project reduce &  Reduce input truth catalogs \\
  &  &   (mag. cut and drop columns) \\ 
do\_01\_build  & rail-project build & Build configurations to run \\
  &  &   truth-to-observed pipeline \\
do\_02\_t2o & rail-project run truth-to-observed & Run truth-to-observed \\
  &  & pipelines to make degraded catalogs \\ 
do\_03\_merge & rail-project merge &  Combine spectroscopic selections \\
do\_04\_subselect & rail-project subsample &  Make train/test files \\
  from catalogs \\
\hline\hline
\end{tabular}
\end{table}



\end{document}

